\section{Experimental Setup}
\label{sec:setup}

\subsection{Benchmark Datasets}

We evaluate the proposed system across two complementary benchmark datasets:

\textbf{Microsoft GUIDE}~\cite{freitas2024cgr} comprises 707,000 verified security
incidents collected from 6,100 customer tenants. Each incident is annotated with
one of three ground-truth labels (TP, BP, FP) and includes structured endpoint
telemetry features (process creation events, network connections, file modifications,
registry accesses). The dataset reflects realistic enterprise alert distributions,
with 20.6\% TP, 49.0\% BP, and 30.4\% FP, reproducing the severe class imbalance
characteristic of production SOC environments.

\textbf{ExCyTIn-Bench}~\cite{wu2026excytin} consists of 57 relational log tables
across 8 complete attack chain scenarios sourced from authentic breach investigations.
Each scenario spans multiple MITRE ATT\&CK phases and requires cross-table join
reasoning to reconstruct the attacker's footstep sequence. The investigation reward
metric $r \in [0, 1]$ measures the fraction of correctly identified attacker steps.

\subsection{Evaluation Protocol}

To prevent cross-tenant data leakage~\cite{arp2022dos}, evaluation employs
5-fold \texttt{GroupKFold} partitioning over Microsoft GUIDE, with groups defined
by customer tenant ID and incident submission timestamp. This ensures that alerts
from the same organizational environment are never simultaneously present in both
training and test splits, preventing optimistic inflation of generalization metrics.

GBDT hyperparameter tuning uses grid search over learning rate $\in \{0.05, 0.1\}$,
maximum tree depth $\in \{4, 6, 8\}$, and subsample fraction $\in \{0.8, 1.0\}$,
with macro-F1 as the selection criterion evaluated on held-out validation tenants.

\subsection{Baseline Comparisons}

We compare against four baselines: (1) Majority Class classifier as a lower bound;
(2) Detector Rule Engine employing fixed threshold logic; (3) Standalone GBDT
(CGR configuration~\cite{freitas2024cgr}) without ReAct planning; and (4) Standalone
LLM agent (AIDA strategy~\cite{tsegai2026aida}) without GBDT priors or Evidence Gate.
